\clearpage
\section{Irreversible Gaussian automata on the chain}\label{sec:chain}

\noindent\motiv{\reve{Chapter~\ref{sec:channels} developed Gaussian channels over an arbitrary real
vector space with an alternating form and extended the definition of a quantum cellular automaton to
completely positive dynamics. We now return to the chain and ask which Gaussian channels commute with
the translations, and how the two notions of locality of Definition~\ref{def:ch-locality} read on
their data. For a channel whose linear map is a rule and whose noise is given by summable cells,
locality becomes a statement about where these two families vanish: finite propagation concerns the
rule alone, while factorization constrains the noise and, through the symplectic form pulled back by
the rule, the rule as well. The thesis asks for factorization because it would let the action on a
single site determine the whole automaton~\cite[p.~91]{krueger} \vetsource{(pinpoint to be
checked)}, and it asks that the class of automata so obtained be closed under concatenation; the
chapter shows how far both wishes are met.}}

\medskip
\noindent\motiv{\reve{Throughout, $G=\Z$, the phase space $V$ and the form $\sigma$ are those of the
lattice, and a rule $L$ has coefficients $\ell_m$, finitely many of them nonzero, with range
$N_L=\{m:\ell_m\neq0\}$ (Corollary~\ref{cor:qca-range}). Cell data $(\nu_x)$ and their covariance
form and symbol are as in Definition~\ref{def:gs-cells}. These notes do not prove the existence of a completely positive map
with prescribed data (Chapter~\ref{sec:channels}). Every statement below
therefore concerns the channels with given data and holds vacuously if there is none, and a
conclusion about a class of channels, such as the failure of closure under concatenation, presupposes
that the channels in question exist; the only channels constructed here are the automorphisms of
symplectic rules.}}

\subsection{Lattice channels}

\begin{definition}[Lattice channel]\label{def:lc-channel}
\reve{A Gaussian channel $T$ on $\CCR(V,\sigma)$ is a \emph{lattice channel} with rule $L$ and noise
cells $(\nu_x)$ if its linear map is $L$, $(\nu_x)$ are cell data whose covariance form is its noise
$\nu$, and its displacement is translation invariant, $\mu(\kappa_{a*}u)=\mu(u)$ for all $a\in\Z$
and $u\in V$.}
\end{definition}

\noindent\motiv{\reve{The automaton $\alpha_L$ of a symplectic rule (Theorem~\ref{thm:rule-qca}) is a
lattice channel with \reve{zero noise cells and zero displacement}, by Proposition~\ref{prop:ch-reversible}. The
running example of the chapter is irreversible.}}

\begin{example}[The lossy chain]\label{ex:lc-lossy}
\reve{Let $n=1$, let $L$ be the example system of Definition~\ref{def:es-example}, fix $0<\zeta<1$, and
consider the data $\zeta L$, with coefficients $\zeta\ell_m$, and the on-site noise cells
$\nu_0=(1-\zeta^2)\cdot1$, $\nu_x=0$ for $x\neq0$. These are the data that
Proposition~\ref{prop:ch-comp} assigns to $\alpha_L\circ T_1$, where $T_1$ has rule $\zeta\cdot1$ and
on-site noise $(1-\zeta^2)\cdot1$, the attenuator of Example~\ref{ex:ch-attenuator} at every site. In
the Schr\"odinger picture this is one step of the example system followed by the loss of the
fraction $1-\zeta^2$ of the intensity at every site into vacuum.}
\end{example}

\noindent\motiv{\reve{Locality is a statement about local algebras, and a channel is known only through
its values on Weyl operators. The following lemma translates between the two.}}

\begin{lemma}[Weyl operators in local algebras]\label{lem:lc-weyl}
\reve{Let $O\subseteq\Z$.
\begin{enumerate}[label=(\alph*)]
\item $\cA(O)$ is the closure of the linear span of the Weyl operators $W(u)$ with $\supp u\subseteq
O$.
\item $W(u)\in\cA(O)$ if and only if $\supp u\subseteq O$.
\end{enumerate}}
\end{lemma}

\begin{proof}\leavevmode
\reve{\begin{enumerate}[label=(\alph*)]
\item The span in question is a $*$-algebra, since a product of Weyl operators is a multiple of the
Weyl operator of the sum, $\supp(u+w)\subseteq\supp u\cup\supp w$ and $W(u)^*=W(-u)$. Its closure is
therefore a closed $*$-algebra containing the generators of $\cA(O)$ and contained in $\cA(O)$, hence
equal to $\cA(O)$.
\item One direction is contained in (a). For the other, suppose that $\supp u\not\subseteq O$. The
Weyl operator $W(v)$ maps the point mass $\delta_0$ in $\ell^2(V)$ to a multiple of modulus one of
the point mass $\delta_{-v}$, so that the continuous linear functional
$\phi(a)=\langle\delta_{-u},a\,\delta_0\rangle$ satisfies $\phi(W(v))\neq0$ exactly for $v=u$. It
vanishes on every generator of $\cA(O)$, and hence on $\cA(O)$ by (a), but not on $W(u)$.
\end{enumerate}}
\end{proof}

\begin{proposition}[Lattice channels are automata]\label{prop:lc-qca}
\reve{A lattice channel $T$ with rule $L$ commutes with every translation, $\alpha_a\circ T=T\circ
\alpha_a$, and $(T,N_L)$ is a quantum cellular automaton.}
\end{proposition}

\begin{proof}
\reve{The channel is unital by Lemma~\ref{lem:ch-unit} and completely positive by definition. Both
$\alpha_a\circ T$ and $T\circ\alpha_a$ are bounded and linear. On a Weyl operator the first
gives $e^{\ii\mu(u)-\nu(u,u)/4}W(\kappa_{a*}Lu)$ and the second
$e^{\ii\mu(\kappa_{a*}u)-\nu(\kappa_{a*}u,\kappa_{a*}u)/4}W(L\kappa_{a*}u)$; these agree, because the
rule commutes with translations, the covariance form of cell data is translation invariant
(Lemma~\ref{lem:gs-cells}(a)) and so is the displacement. Lemma~\ref{lem:ch-unit} gives equality.
For locality, $T(W(u))$ is a multiple of $W(Lu)$ and $\supp Lu\subseteq\supp u+N_L$ by
Corollary~\ref{cor:qca-range}, so that $T$ maps the Weyl operators of $\cA(O)$ into the closed
subalgebra $\cA(O+N_L)$; by continuity and Lemma~\ref{lem:lc-weyl}(a) it maps $\cA(O)$ there.}
\end{proof}

\subsection{Pull-back cells and locality}

\noindent\motiv{\reve{Factorization involves the product formula of
Proposition~\ref{prop:ch-product}, and hence the dissipation form
$\nu+\ii(\sigma-\sigma(L\cdot,L\cdot))$. Its three terms are translation-invariant forms, and we
describe each by its cells. For the noise these are the given cells; for $\sigma$ they are $\Omega$ at
the origin and zero elsewhere; and for $\sigma(L\cdot,L\cdot)$ they are obtained, as in the evolution
of a covariance in Theorem~\ref{thm:gs-evolution}, by pulling the cells of $\sigma$ back along the
rule.}}

\begin{definition}[Pull-back cells]\label{def:lc-pullback}
\reve{For a rule $L$ and a family $(\nu_x)$ of real $2n\times2n$ matrices, the \emph{pull-back cells}
are
\[(L_*\nu)_j=\sum_{m,m'}\ell_m^{\mathsf T}\,\nu_{j+m-m'}\,\ell_{m'}\qquad(j\in\Z),\]
the sum running over the finitely many $m,m'$ with $\ell_m,\ell_{m'}\neq0$. We write
$(L_*\sigma)_j=\sum_m\ell_m^{\mathsf T}\Omega\,\ell_{m+j}$ for the pull-back cells of the on-site
cells $\Omega\delta_{x,0}$ of $\sigma$, and $\delta_{\bullet,0}\,\nu_0$ for the on-site cells equal to a
matrix $\nu_0$ at the origin and to zero elsewhere. These are the cells $L_*\gamma$ of
Chapter~\ref{sec:stationary}, written there for cell data; the subscript records that forms are
pulled back along $L$, an operation often written $L^*$ in geometry. The \emph{form} of an arbitrary family
$(\nu_x)$ is $\nu(u,w)=\sum_{x,y}u(x)^{\mathsf T}\nu_{x-y}w(y)$, a finite sum; it determines the family,
since $\nu(e_{x,i},e_{0,j})=(\nu_x)_{ij}$.}
\end{definition}

\begin{lemma}[Forms of pull-back cells]\label{lem:lc-pullback}
\reve{Let $L$ be a rule.
\begin{enumerate}[label=(\alph*)]
\item For every family $(\nu_x)$, the form $\sum_{x,y}u(x)^{\mathsf T}(L_*\nu)_{x-y}w(y)$ is
$\nu(Lu,Lw)$, where $\nu$ is the form of the cells $(\nu_x)$; if $(\nu_x)$ are cell data, so are
$(L_*\nu)$.
\item $\sigma(Lu,Lw)=\sum_{x,y}u(x)^{\mathsf T}(L_*\sigma)_{x-y}w(y)$, and $(L_*\sigma)_j$ is the
coefficient of $z^j$ in $\hat L^{\star\mathsf T}\Omega\hat L$, where $\hat L^\star$ replaces $z$ by
$z^{-1}$.
\item If $(\nu_x)$ are cell data, the symbol of $L_*\nu$ is
$\widehat{L_*\nu}(k)=\hat L(k)^\dagger\,\hat\nu(k)\,\hat L(k)$.
\end{enumerate}}
\end{lemma}

\begin{proof}\leavevmode
\reve{\begin{enumerate}[label=(\alph*)]
\item The cells $(L_*\nu)$ are the cells $\gamma'$ of Theorem~\ref{thm:gs-evolution}, and the
computation of the form in its proof uses neither summability nor symmetry of the cells; the second
statement is part of that theorem.
\item The first statement is (a) for the on-site cells of $\sigma$, whose form is $\sigma$. In
$\hat L^{\star\mathsf T}\Omega\hat L=\sum_{m,m'}\ell_m^{\mathsf T}\Omega\,\ell_{m'}z^{m'-m}$ the
coefficient of $z^j$ collects the terms with $m'=m+j$.
\item This is the formula for the symbol in Theorem~\ref{thm:gs-evolution}.
\end{enumerate}}
\end{proof}

\begin{theorem}[Locality on the chain]\label{thm:lc-locality}
\reve{Let $T$ be a lattice channel with rule $L$ and noise cells $(\nu_x)$.
\begin{enumerate}[label=(\Roman*)]
\item For a finite set $N\subseteq\Z$, $T$ has finite propagation with neighbourhood $N$ if and only
if $N_L\subseteq N$.
\item For a finite set $N'\subseteq\Z$ with $0\in N'=-N'$, $T$ factorizes with respect to $N'$ if
and only if $\nu_x=0$ and $(L_*\sigma)_x=0$ for every $x\notin N'$.
\end{enumerate}}
\end{theorem}

\begin{proof}\leavevmode
\reve{\begin{enumerate}[label=(\Roman*)]
\item If $N_L\subseteq N$, then $T$ has finite propagation with neighbourhood $N_L$ by
Proposition~\ref{prop:lc-qca}, and hence with $N$. Conversely, for the site vector $u=e_{0,i}$,
$T(W(u))$ is a nonzero multiple of $W(Lu)$ and lies in $\cA(N)$, so that $\supp Lu\subseteq N$ by
Lemma~\ref{lem:lc-weyl}(b). Since $(Lu)(m)$ is the $i$-th column of $\ell_m$, every $m$ for which this
column is nonzero lies in $N$, and taking all $i$ gives $N_L\subseteq N$.
\item Suppose first that the cells vanish outside $N'$, and let $u,w$ have supports with $(\supp
u+N')\cap\supp w=\emptyset$, so that $x-y\notin N'$ for $x\in\supp u$ and $y\in\supp w$. By
Lemma~\ref{lem:lc-pullback}, $\nu(u,w)=0$ and $\sigma(Lu,Lw)=0$, and $\sigma(u,w)=0$ because the
supports are disjoint; the dissipation form vanishes on $(u,w)$, and
$T(W(u)W(w))=T(W(u))T(W(w))$ by Proposition~\ref{prop:ch-product}. For regions $O_1,O_2$ with
$(O_1+N')\cap O_2=\emptyset$, the maps $(b_1,b_2)\mapsto T(b_1b_2)$ and $(b_1,b_2)\mapsto
T(b_1)T(b_2)$ are bilinear and continuous on $\cA(O_1)\times\cA(O_2)$ and agree on pairs of Weyl
operators, hence everywhere by Lemma~\ref{lem:lc-weyl}(a). Conversely, suppose that $T$ factorizes with respect to $N'$, let $x\notin N'$, and let $u=e_{x,i}$ and
$w=e_{0,j}$ be site vectors. The regions $\{x\}$ and $\{0\}$
are separated by $N'$, and for every real $t\neq0$ factorization together with the product formula
gives $e^{-\Delta_T(tu,w)/2}\,T(W(tu))\,T(W(w))=T(W(tu))\,T(W(w))$, where the product on the right is
a nonzero multiple of a unitary. Hence $e^{-t\Delta_T(u,w)/2}=1$ for every $t\neq0$; choosing $t$ so
small that $|t\Delta_T(u,w)/2|<2\pi$ forces $\Delta_T(u,w)=0$. Its real part is
$\nu(u,w)=(\nu_x)_{ij}$, and its imaginary part is $\sigma(u,w)-\sigma(Lu,Lw)=-((L_*\sigma)_x)_{ij}$,
because $\sigma(u,w)=0$ for $x\neq0$.
\end{enumerate}}
\end{proof}

\noindent\motiv{\reve{For $N'=\{0\}$ the condition on the rule is that
$(L_*\sigma)_x=\sum_m\ell_m^{\mathsf T}\Omega\,\ell_{m+x}$ vanish for $x\neq0$. These are the coefficient
equations~\eqref{eq:qca-coeff} of a symplectic rule \reve{for $\Omega$ at every index other than $0$,} without the
\reve{equation at the index $0$}; the rule need not be symplectic, and $\hat L^{\star\mathsf T}\Omega\hat L$ need only
be a constant matrix. The two notions of locality are independent of each other in the following
sense.}}

\begin{proposition}[Independence of the two notions]\label{prop:lc-independent}
\reve{Let $n\ge1$.
\begin{enumerate}[label=(\alph*)]
\item There is a Gaussian channel that factorizes with respect to $\{0\}$ and does not have finite
propagation with neighbourhood $\{0\}$.
\item For every finite set $N'\subseteq\Z$ with $0\in N'=-N'$ there are cell data $(\nu_x)$ of
finite range such that the rule $0$ and the noise $\nu$ satisfy~\eqref{eq:ch-cond}, and no lattice
channel with these data factorizes with respect to $N'$.
\end{enumerate}}
\end{proposition}

\begin{proof}\leavevmode
\reve{\begin{enumerate}[label=(\alph*)]
\item The shift, the rule with coefficient $\ell_1=1$ and all others zero, is symplectic, and its
automorphism is a lattice channel. Since $\hat L^{\star\mathsf T}\Omega\hat L=z^{-1}\Omega z=\Omega$,
the pull-back cells of $\sigma$ vanish off the origin, and Theorem~\ref{thm:lc-locality}(II) gives
factorization with respect to $\{0\}$; since $N_L=\{1\}$, part (I) excludes finite propagation with
neighbourhood $\{0\}$.
\item Let $x_0=1+\max_{m\in N'}|m|$, and let $\nu_0=2\cdot1$, $\nu_{\pm x_0}=\frac12\cdot1$ and
$\nu_x=0$ otherwise. For the rule $0$, condition~\eqref{eq:ch-cond} is the state condition
$\nu+\ii\sigma\ge0$ (Proposition~\ref{prop:ch-special}(c)), which holds by
Theorem~\ref{thm:gs-fourier}(a), since
$\hat\nu(k)+\ii\Omega=(1+\cos kx_0)\cdot1+(1+\ii\Omega)$ is a sum of positive semidefinite
matrices. Since $x_0\notin N'$ and $\nu_{x_0}\neq0$, Theorem~\ref{thm:lc-locality}(II) excludes
factorization with respect to $N'$ for every lattice channel with these data.
\end{enumerate}}
\end{proof}

\noindent\motiv{\reve{Every lattice channel with the rule $0$ has finite propagation with neighbourhood
$\{0\}$ by Theorem~\ref{thm:lc-locality}(I), since \reve{$N_L=\emptyset$ for $L=0$}. Part (b) therefore shows that
finite propagation does not imply factorization, and part (a) that factorization does not imply
finite propagation with a prescribed neighbourhood.}}

\subsection{Concatenation}

\noindent\motiv{\reve{Applying two lattice channels in succession gives a lattice channel whose rule is
the product of the two rules and whose noise collects the first noise and the pull-back of the second.
The ranges of the factors control those of the resulting families.}}

\begin{proposition}[Concatenation of lattice channels]\label{prop:lc-comp}
\reve{If $T_1$ and $T_2$ are lattice channels with rules $L_1,L_2$ and noise cells $(\nu_{1,x})$,
$(\nu_{2,x})$, then $T_2\circ T_1$ is a lattice channel with rule $L_2L_1$ and noise cells
$\nu_1+L_{1*}\nu_2$.}
\end{proposition}

\begin{proof}
\reve{By Proposition~\ref{prop:ch-comp}, $T_2\circ T_1$ is a Gaussian channel with linear map
$L_2L_1$, noise $\nu_1+\nu_2(L_1\cdot,L_1\cdot)$ and displacement $\mu_1+\mu_2\circ L_1$. The noise is
the form of the cells $\nu_1+L_{1*}\nu_2$ by Lemma~\ref{lem:lc-pullback}(a), which are cell data, and
the displacement is translation invariant because $L_1$ commutes with translations.}
\end{proof}

\begin{proposition}[Ranges under concatenation]\label{prop:lc-supports}
\reve{Let $L,L_1,L_2$ be rules.
\begin{enumerate}[label=(\alph*)]
\item $N_{L_2L_1}\subseteq N_{L_1}+N_{L_2}$.
\item If $(\nu_x)$ vanishes outside a finite set $M$, then $(L_*\nu)_j=0$ for $j\notin
M+(N_L-N_L)$.
\item $(L_*\sigma)_j=0$ for $j\notin N_L-N_L$.
\item Let $N_1,N_2$ be finite, $N_2$ nonempty, $N_{L_i}\subseteq N_i$, and let $(\nu_{i,x})$ vanish
outside $N_i-N_i$ for $i=1,2$. Then $N_{L_2L_1}\subseteq N_1+N_2$, and $\nu_1+L_{1*}\nu_2$ vanishes
outside $(N_1+N_2)-(N_1+N_2)$.
\item For on-site noises $\nu_{i,x}=\delta_{x,0}\nu_{i,0}$, the noise cells $\nu_1+L_{1*}\nu_2$ vanish
off the origin if and only if $(L_{1*}\nu_2)_x=0$ for every $x\neq0$.
\end{enumerate}}
\end{proposition}

\begin{proof}\leavevmode
\reve{\begin{enumerate}[label=(\alph*)]
\item The coefficient of $z^j$ in $\hat L_2\hat L_1$ is a sum of products of coefficients of $z^b$ in
$\hat L_2$ and of $z^a$ in $\hat L_1$ with $a+b=j$, and is zero unless such $a\in N_{L_1}$ and $b\in
N_{L_2}$ exist.
\item A term of $(L_*\nu)_j$ is nonzero only if $j+m-m'\in M$ for some $m,m'\in N_L$, that is,
$j\in M+(N_L-N_L)$.
\item This is (b) for the on-site cells of $\sigma$, with $M=\{0\}$.
\item The first statement follows from (a). Choose $c\in N_2$. If $x=a-b$ with $a,b\in N_1$, then
$x=(a+c)-(b+c)\in(N_1+N_2)-(N_1+N_2)$, so that $\nu_1$ vanishes outside this set. By (b), the cells
$L_{1*}\nu_2$ vanish outside $(N_2-N_2)+(N_{L_1}-N_{L_1})$, and an element $(a-b)+(a'-b')$ of this
set, with $a,b\in N_2$ and $a',b'\in N_{L_1}\subseteq N_1$, is $(a'+a)-(b'+b)$.
\item The first summand vanishes off the origin.
\end{enumerate}}
\end{proof}

\noindent\motiv{\reve{The inclusion in (a) can be strict, since coefficients may cancel: the rules
with symbols $(1-z)\cdot1$ and $(1+z)\cdot1$ have ranges $\{0,1\}$, and their product $(1-z^2)\cdot1$
has range $\{0,2\}$, a proper subset of $\{0,1,2\}$.}}

\noindent\motiv{\reve{Part (d) says that channels whose noise is correlated only within the
differences of their neighbourhoods form a class closed under concatenation. The channels that
factorize with respect to $\{0\}$, whose noise and whose pull-back of $\sigma$ are both on site, do
not, as the next \reve{subsection} shows for the noise; part (e) says what closure would require of it.}}

\noindent\motiv{\reve{The pull-back of $\sigma$ can also leave the site under concatenation. Since a
family is determined by its form, Lemma~\ref{lem:lc-pullback}(a) applied twice gives
$(L_2L_1)_*\nu=L_{1*}(L_{2*}\nu)$ for every family. If $L_{1*}\sigma$ and $L_{2*}\sigma$ are on site
and $n=1$, the on-site cell of $L_{2*}\sigma$ is antisymmetric and hence a multiple $c\,\Omega$, and
then $(L_2L_1)_*\sigma=c\,L_{1*}\sigma$ is again on site. For $n=2$ this fails. With coordinates
$(q_1,q_2,p_1,p_2)$, let $L_1$ have the coefficients $\ell_0=P_2$ and $\ell_1=P_1$, the projections
onto the second and the first mode, so that it shifts the first mode and fixes the second, and let
$L_2$ act at every site by the matrix $S=1+E_{12}$, that is, $q_1\mapsto q_1+q_2$. Both $L_{1*}\sigma$
and $L_{2*}\sigma$ are on site, the latter with the cell $S^{\mathsf T}\Omega S=\Omega+E_{23}-E_{32}$,
but the cell of $(L_2L_1)_*\sigma$ at $j=1$ is $\ell_0^{\mathsf T}(S^{\mathsf T}\Omega S)\ell_1=
P_2(E_{23}-E_{32})P_1=E_{23}\neq0$.}}

\subsection{The example system}

\noindent\motiv{\reve{For the example system, we can compute the cell at distance two of the pull-back of on-site
noise in closed form, and this cell decides when two steps with on-site noise again have
on-site noise.}}

\begin{theorem}[Two steps with on-site noise]\label{thm:lc-twosteps}
\reve{Let $n=1$, let $L$ be the example system of Definition~\ref{def:es-example}, with
$r=(\sin\varphi,\cos\varphi)^{\mathsf T}$ the transposed second row of $\ell_0$, and let $\nu_0$ be a
real $2\times2$ matrix. Then
\[(L_*\nu)_2=(\nu_0)_{11}\,\lambda^2\,rr^{\mathsf T}\]
for the on-site cells $\nu_x=\delta_{x,0}\nu_0$. If $\nu_0$ is symmetric and positive semidefinite,
then
$(L_*\nu)_x=0$ for every $x\neq0$ if and only if $\lambda=0$ or $(\nu_0)_{11}=0$.}
\end{theorem}

\begin{proof}
\reve{For on-site cells, $(L_*\nu)_j=\sum_{m'-m=j}\ell_m^{\mathsf T}\nu_0\ell_{m'}$. For $j=2$ only
$(m,m')=(-1,1)$ occurs, and with $\ell_{\pm1}=\lambda e_1r^{\mathsf T}$,
\[\ell_{-1}^{\mathsf T}\nu_0\ell_1=\lambda^2\,r\,(e_1^{\mathsf T}\nu_0e_1)\,r^{\mathsf
T}=\lambda^2(\nu_0)_{11}\,rr^{\mathsf T}.\]
If the cells vanish off the origin, the trace of this cell, $\lambda^2(\nu_0)_{11}|r|^2=
\lambda^2(\nu_0)_{11}$, vanishes. Conversely, if $\lambda=0$ only $m=m'=0$ contributes. If
$(\nu_0)_{11}=0$ and $\nu_0$ is symmetric and positive semidefinite, then $\nu_0e_1=0$, since such a
matrix with a zero diagonal entry has a zero row and column; hence
$\nu_0\ell_{\pm1}=0=\ell_{\pm1}^{\mathsf T}\nu_0$, and every term with $j\neq0$ contains some
$\ell_{\pm1}$.}
\end{proof}

\noindent\motiv{\reve{By the congruence $\gamma_B=\Omega^{\mathsf T}\gamma\,\Omega$ of
Chapter~\ref{sec:gaussianstates}, which exchanges the diagonal entries of a $2\times2$ form, the
entry $(\nu_0)_{11}$ of the noise form is the momentum entry of the matrix that the noise adds to the
covariance matrix of each site, that is, twice the variance it adds to the momentum. Pulled back by
one step of the rule, this part of the noise correlates sites two apart through the two coupling
factors in $\ell_{-1}^{\mathsf T}\nu_0\ell_1$, while noise with $(\nu_0)_{11}=0$, such as
$\nu_0=\operatorname{diag}(0,b)$, which adds to the position alone, is not read by the coupling
$\ell_{\pm1}=\lambda e_1r^{\mathsf T}$ and stays on site. For the lossy chain the consequence is the following.}}

\begin{proposition}[The lossy chain]\label{prop:lc-lossy}
\reve{Let $0<\zeta<1$, and let $\zeta L$ and the on-site cells $(1-\zeta^2)\cdot1$ be the data of
Example~\ref{ex:lc-lossy}.
\begin{enumerate}[label=(\alph*)]
\item The data satisfy~\eqref{eq:ch-cond}.
\item $((\zeta L)_*\sigma)_j=0$ for every $j\neq0$.
\item The rule $\zeta L$ is not symplectic.
\item If $\lambda\neq0$, the noise cells of two steps, $\nu+(\zeta L)_*\nu$, are nonzero at $j=2$.
\end{enumerate}}
\end{proposition}

\begin{proof}\leavevmode
\reve{\begin{enumerate}[label=(\alph*)]
\item Since $L$ is symplectic, $\sigma(\zeta Lu,\zeta Lw)=\zeta^2\sigma(u,w)$, and the matrices
of~\eqref{eq:ch-cond} are $(1-\zeta^2)$ times those of $\gamma+\ii\sigma$, with $\gamma$ the vacuum form
of the cells $\delta_{x,0}\cdot1$; the vacuum form satisfies the state condition
(Example~\ref{ex:gs-vacuum}).
\item By Lemma~\ref{lem:lc-pullback}(b), $((\zeta L)_*\sigma)_j$ is the coefficient of $z^j$ in
$\zeta^2\hat L^{\star\mathsf T}\Omega\hat L=\zeta^2\Omega$, since the example system is symplectic
(Example~\ref{ex:qca-rotation}).
\item $(\zeta\hat L)^{\star\mathsf T}\Omega(\zeta\hat L)=\zeta^2\Omega\neq\Omega$, so that $\zeta L$ is
not symplectic by Theorem~\ref{thm:qca-criterion}.
\item The on-site first summand does not contribute at $j=2$. The pull-back cells of the rule $\zeta L$
are $\zeta^2$ times those of $L$, so that by Theorem~\ref{thm:lc-twosteps} with
$\nu_0=(1-\zeta^2)\cdot1$ the second summand at $j=2$ is $\zeta^2(1-\zeta^2)\lambda^2rr^{\mathsf
T}\neq0$.
\end{enumerate}}
\end{proof}

\noindent\motiv{\reve{By Theorem~\ref{thm:lc-locality}, every lattice channel with the data of the lossy
chain therefore has finite propagation with neighbourhood $\{-1,0,1\}$ and factorizes with respect to
$\{0\}$, although its rule is not symplectic, whereas two steps of it, a lattice channel by
Proposition~\ref{prop:lc-comp}, no longer factorize with respect to $\{0\}$ when $\lambda\neq0$.}}

\subsection{The Fourier form of the condition and dilations}

\noindent\motiv{\reve{For translation-invariant data, condition~\eqref{eq:ch-cond} reduces, as the
state condition does in Theorem~\ref{thm:gs-fourier}, to one Hermitian matrix at each wave number.}}

\begin{theorem}[Fourier form of the condition]\label{thm:lc-fourier}
\reve{Let $(\nu_x)$ be cell data of finite range and $L$ a rule. Then $L$ and $\nu$
satisfy~\eqref{eq:ch-cond} if and only if
\[\hat\nu(k)+\ii\Omega-\ii\,\hat L(k)^\dagger\Omega\hat L(k)\ge0\qquad\text{for every }k\in\R.\]}
\end{theorem}

\begin{proof}
\reve{By Lemma~\ref{lem:lc-pullback}, the Hermitian form of~\eqref{eq:ch-cond} on complex combinations
of vectors, $h(f)=\sum_{x,y}f(x)^\dagger a_{x-y}f(y)$ for finitely supported
$f\colon\Z\to\C^{2n}$, has the cells $a_j=\nu_j+\ii(\Omega\delta_{j,0}-(L_*\sigma)_j)$, of finite
range, with $a_{-j}=a_j^\dagger$, and
$\sum_ja_je^{-\ii kj}=\hat\nu(k)+\ii\Omega-\ii\hat L(k)^\dagger\Omega\hat L(k)$ by
Lemma~\ref{lem:lc-pullback}(b) and the computation of the symbol in Theorem~\ref{thm:gs-evolution}.
The proof of Lemma~\ref{lem:gs-parseval} and Theorem~\ref{thm:gs-fourier} uses of the cells
$a_g=\gamma_g+\ii\delta_{g,0}\Omega$ only that they are summable and satisfy $a_{-g}=a_g^\dagger$,
and the necessity part only their finite range; it applies verbatim to the cells $a_j$ and gives
both directions.}
\end{proof}

\noindent\motiv{\reve{We assume finite range only for simplicity, since the cells $a_j$ of summable
noise are summable and Remark~\ref{rem:gs-fourier-general} extends the necessity to them.} \reve{At
each wave number, the matrix of the theorem is the matrix $\nu+\ii(\Omega-L^{\mathsf T}\Omega L)$ whose
positivity is the condition for finitely many modes (proof of Proposition~\ref{prop:ch-matrices}(c)),
with the real matrices $L$ and $\nu$ replaced by the complex matrices $\hat L(k)$ and $\hat\nu(k)$ and
the transpose by the conjugate transpose, so that the condition on the chain is the finite-mode
condition of Chapter~\ref{sec:channels} read at every wave number.}}

\noindent\motiv{\reve{For the lossy chain, $\hat\nu(k)=(1-\zeta^2)\cdot1$ and
$(\zeta\hat L(k))^\dagger\Omega(\zeta\hat L(k))=\zeta^2\Omega$, since the example system is symplectic,
so that the matrix of the theorem is $(1-\zeta^2)(1+\ii\Omega)$ at every wave
number, which is positive semidefinite, in agreement with Proposition~\ref{prop:lc-lossy}(a).}}

\noindent\motiv{\reve{It remains to consider dilations on the chain. Let a translation-invariant ancilla with $n'$ modes per site be coupled to the
chain by a rule $S$ on $n+n'$ modes per site that preserves the form of
$\Omega_n\oplus\Omega_{n'}$. As in Proposition~\ref{prop:ch-dilation}, $S(u,0)=(S_{11}u,S_{21}u)$,
where the symbol $\hat S_{11}$ of the system block is a $2n\times2n$ and the symbol $\hat S_{21}$ of
the block from system to ancilla a $2n'\times2n$ matrix of Laurent polynomials; we define the pull-back cells
$S_{21*}\gamma$ of cells of the ancilla as in Definition~\ref{def:lc-pullback}.}}

\begin{theorem}[Dilations on the chain]\label{thm:lc-dilation}
\reve{Let $\hat S^{\star\mathsf T}(\Omega_n\oplus\Omega_{n'})\hat S=\Omega_n\oplus\Omega_{n'}$, and
let $\gamma_0$ be a real $2n'\times2n'$ matrix with $\gamma_0+\ii\Omega_{n'}\ge0$, which forces
$\gamma_0$ to be symmetric. Then the rule
$S_{11}$ and the noise with cells $\nu=S_{21*}(\delta_{\bullet,0}\gamma_0)$
satisfy~\eqref{eq:ch-cond}; $\hat\nu(k)=\hat S_{21}(k)^\dagger\gamma_0\hat S_{21}(k)$ for every $k$;
$N_{S_{11}}\subseteq N_S$; and $\nu_j=0$ for $j\notin N_S-N_S$.}
\end{theorem}

\begin{proof}
\reve{The block of the system in $\hat S^{\star\mathsf T}(\Omega_n\oplus\Omega_{n'})\hat S=\Omega_n
\oplus\Omega_{n'}$ reads $\hat S_{11}^{\star\mathsf T}\Omega_n\hat S_{11}+\hat
S_{21}^{\star\mathsf T}\Omega_{n'}\hat S_{21}=\Omega_n$, which by the computation of
Lemma~\ref{lem:lc-pullback}(b) says that $\sigma(S_{11}u,S_{11}w)+\sigma'(S_{21}u,S_{21}w)=
\sigma(u,w)$, with $\sigma'$ the form of the ancilla; this is
Proposition~\ref{prop:ch-dilation}(a) for the translation-invariant maps $S_{11}$ and $S_{21}$, and
Proposition~\ref{prop:ch-dilation}(b) gives~\eqref{eq:ch-cond}, since the on-site form of
$\gamma_0+\ii\Omega_{n'}$ is positive semidefinite. We compute the symbol as in
Theorem~\ref{thm:gs-evolution} for the rectangular matrix $\hat S_{21}$, and the ranges follow as in
Proposition~\ref{prop:lc-supports}\reve{(b)}, because the blocks of $\hat S$ have ranges within
$N_S$.}
\end{proof}

\noindent\motiv{\reve{For the lossy chain, let $B$ be the beam splitter of
Example~\ref{ex:ch-beamsplitter} at every site, with one ancilla mode per site, and let
$S=(L\oplus1)\circ B$, which preserves $\Omega_1\oplus\Omega_1$ since both factors do. On a system
vector, $S(u,0)=(\zeta Lu,-\sqrt{1-\zeta^2}\,u)$, so that $S_{11}=\zeta L$ and
$S_{21}=-\sqrt{1-\zeta^2}\cdot1$,
and with the ancilla in the vacuum, $\gamma_0=1$, the theorem gives the on-site noise
$(1-\zeta^2)\cdot1$ of the lossy chain. The composite $\alpha_{L\oplus1}\circ\alpha_B$ acts on states
first by $\alpha_{L\oplus1}$, so that this is one step of the example system followed by the beam
splitter, as in Example~\ref{ex:lc-lossy}. The other order, $S=B\circ(L\oplus1)$, gives
$S_{21}=-\sqrt{1-\zeta^2}\,L$ and the noise cells $(1-\zeta^2)L_*(\delta_{\bullet,0}\cdot1)$, which by the
computation of Theorem~\ref{thm:lc-twosteps} are nonzero at $j=2$ when $\lambda\neq0$.}}

\subsection{Notes and references}\label{sec:lc-notes}

\begin{enumerate}[label=(\arabic*)]
\item \reve{\emph{Lattice channels.} Definition~\ref{def:lc-channel} is the translation-invariant
setting of~\cite[p.~92]{krueger}, whose rule and noise are $\Gamma(x)$ and $g(x)$, the rule there of
finite range. The notes ask for summable noise cells, as in Chapter~\ref{sec:gaussianstates};
Theorem~\ref{thm:lc-fourier} assumes finite range, which the data of
Proposition~\ref{prop:lc-independent}(b) have.}
\item \reve{\emph{Locality.} Theorem~\ref{thm:lc-locality} is Lemma~4.11 of~\cite[p.~95]{krueger},
which is stated with ``if''; its proof argues the converse for notion (II), and the scaling step
that turns $e^{-\Delta_T/2}=1$ into $\Delta_T=0$ is not written there. The thesis's condition
(4.36), $\sum_x\Gamma^+_x\Gamma_{\Delta+x}=0$ with $\Gamma^+=-\sigma_s\Gamma^{\mathsf T}\sigma_s$, is
$-\Omega(L_*\sigma)_\Delta=0$. The thesis's witness for ``(II) does not imply (I)'' in the proof of
Lemma~4.10(v)~\cite[pp.~94--95]{krueger}, the completely depolarizing channel, maps every observable
to a multiple of the unit and so has finite propagation with neighbourhood $\{0\}$. Since every
automaton has finite propagation for some finite neighbourhood, what can be meant is that
factorization gives no bound on the neighbourhood, which Proposition~\ref{prop:lc-independent}(a)
shows for the neighbourhood $\{0\}$ and the same argument with the shift by $1+\max_{m\in N}|m|$ for
any \reve{nonempty} finite $N$. Part (b) makes explicit the thesis's witness for the converse, noise of range beyond
the given set.}
\item \reve{\emph{Concatenation.} Proposition~\ref{prop:lc-comp} and
Proposition~\ref{prop:lc-supports} are the discussion of~\cite[p.~96]{krueger}, where the ranges
of the composite are written as $N_1+N_2$ and $M_1\cup(M_2+N_1-N_1)$; only the inclusions hold, by the
example after the proof for the first and by Theorem~\ref{thm:lc-twosteps} for the second, where the
noise $\operatorname{diag}(0,b)$ has an on-site pull-back although $M_2+N_L-N_L=\{-2,\dots,2\}$. The
nonempty $N_2$ in (d) is supplied in the thesis by its convention $0\in N$. Part (e) is (4.39).}
\item \reve{\emph{The example system.} The thesis states that the noise of two steps does not fulfil
(4.39) whenever the coupling is nonzero~\cite[p.~97]{krueger}; by Theorem~\ref{thm:lc-twosteps} this
holds for isotropic noise $c\cdot1$, $c>0$, and fails for noise with $(\nu_0)_{11}=0$, which adds
variance to the position quadrature alone, such as $\nu_0=\operatorname{diag}(0,b)$, $b>0$, with the
symplectic rule $L$, for which these data satisfy~\eqref{eq:ch-cond} by
Proposition~\ref{prop:ch-special}(b). The remark after Lemma~4.11 of~\cite[p.~95]{krueger}
asserts that for one mode per site a rule of both types need not be symplectic; the lossy chain is an
explicit instance, with noise that satisfies~\eqref{eq:ch-cond}.}
\item \reve{\emph{Dilations.} Theorem~\ref{thm:lc-dilation} is Lemma~4.12 of~\cite[p.~97]{krueger} at
the level of covariances, with the reduction (4.40) of~\cite[p.~98]{krueger}; that the data satisfy
the condition goes beyond the thesis, which checks that the noise is real and symmetric. The thesis
writes the symbol of the noise with transposes, $\hat C^{\mathsf T}(k)\hat\gamma'\hat C(k)$; the
conjugate transpose of Theorem~\ref{thm:lc-dilation} agrees with it for reflection-symmetric blocks,
as in the note on (4.17) in Chapter~\ref{sec:gaussianstates}. The range statements correspond to
Lemma~4.10(iii)--(iv) of~\cite[p.~94]{krueger}. The construction of the channel on the algebra is not
carried out here. In finite dimensions without translation invariance, factorization does not imply
that a channel is a local dilation~\cite{piroli-cirac} \vetsource{(pinpoint to be checked)}.}
\item \reve{\emph{Not treated here.} Section~4.3 of the thesis also defines two further notions of
locality, (III) by local Kraus operators and (IV) by local dilations, relates the four notions in
Lemma~4.10~\cite[pp.~93--94]{krueger}, and leaves open whether every irreversible Gaussian automaton
is of type (IV)~\cite[p.~98]{krueger} \vetsource{(pinpoint to be checked)}.}
\end{enumerate}
